\chapter{Carte de l'ouvrage et methode de construction}

Le livre suivra les trois grands ensembles de la specialisation : apprentissage
supervise, algorithmes avances, puis apprentissage non supervise, recommandation
et apprentissage par renforcement. Chaque notion sera reliee aux fils
transverses utiles a l'ingenierie : biometrie, identite, qualite, risque,
observabilite et apprentissage continu.

\section{Un chapitre en sept couches}

Chaque chapitre vise sept niveaux de comprehension :
\begin{enumerate}
  \item le probleme concret;
  \item l'intuition visuelle;
  \item la notation et sa lecture orale;
  \item la derivation mathematique;
  \item l'implementation reproductible;
  \item les limites et les risques d'interpretation;
  \item le transfert vers les projets d'ingenierie.
\end{enumerate}

\section{Prochaines integrations}

Les prochaines versions consolideront d'abord les notes deja disponibles sur
K-means, les arbres de decision et les ensembles. Les nouveaux chapitres seront
ensuite ajoutes au fil de la progression, avec une fiche de revision et un TP
pour chaque bloc conceptuel important.

